\section{Discussion}
\label{discussion:section}

The phase framework separates two tasks: a diagonal unitary supplies a complete orthogonal measurement, while the interference terms determine its local geometry. A fixed four-feature Boolean kernel makes these local terms coincide for every $n\ge5$, producing a continuous radius interval from $0$ to $2r_E(n)$. The classification of radius laws for Boolean phase lines shows that this endpoint is optimal within the full two-valued-generator model. For even $n\ge8$, cancellation of the mixed background terms places the known cosine family and the enhanced family in one genuinely two-dimensional measurement family. The results at $n=3$ and $n=4$ explain the lower-dimensional boundary of this picture: exact local rigidity at the specified point at $n=3$, and a distinction between straight and curved phase deformations at $n=4$. Throughout these constructions, the actual single-qubit reduced states provide the link between the algebraic phase conditions, fixed tetrahedral directions, and continuously variable single-qubit-versus-rest entanglement.

Two questions remain close to these results. At $n=4$, the specified straight segment of the larger-radius family in Ref.~\cite{PauwelsGisin2026} contains no measurement equivalent to our reference EJM under local unitaries and relabelings. Whether the curved branch constructed here connects to that family elsewhere in the regular tetrahedral phase space is unresolved. At $n=6$, the four-color model has only the two crossing branches near the reference point. A finer phase partition may allow a two-parameter family containing both branches; resolving this requires controlling additional mixed terms beyond those present in the model studied here.
